% ============================================================================
% Chapter 3 --- The sovereignty grid (full version)
% ----------------------------------------------------------------------------
% Purpose : publish the normative pivot of the entire document. The grid is
%           the criterion that determines whether a candidate component is
%           acceptable; it transforms the OSS-vs-proprietary debate into a
%           multi-criteria arbitration.
% Page target: ~8 pages.
% Dependencies: foundational, referenced by every Part-II and Part-III chapter.
% ----------------------------------------------------------------------------
% Decisions registered in _coherence-EN.md:
%   D3.1 five orthogonal dimensions, never combined into a single score
%   D3.2 fixed four-point scoring scale (0/1/2/3), no half-points
%   D3.3 deliberation instrument, not audit instrument
%   D3.4 data dimension first by convention
% ============================================================================

\chapter{The sovereignty grid}
\label{ch:sovereignty-grid}

Sovereignty in university \gls{it} is rarely contested in principle and
rarely operationalised in practice. The grid presented in this chapter
is the instrument that bridges the two: a structured way of asking, for
any candidate component of the digital infrastructure, whether the
institution retains the capacities --- legal, technical, financial,
operational, and institutional --- that allow it to remain the author
of its own decisions. The grid is the normative pivot of this Technical
Reference: every subsequent chapter applies it; every recommendation
derives from it.

\begin{caveatbox}[Posture of this chapter]
\noindent The grid is a deliberation and procurement instrument, not an
audit scoring system. A component admissible on the grid is not, by
that fact alone, the right choice; the institutional context, which the
grid does not capture in full, completes the decision. Conversely, a
component inadmissible on a critical dimension cannot be rescued by
strengths elsewhere: each dimension is reported independently and
critical thresholds may apply per workload class.
\end{caveatbox}

%-----------------------------------------------------------------------------
\section{Five dimensions of sovereignty}
\label{sec:sg-dimensions}

Sovereignty is too often treated as a single concept. In practice, an
institution can be sovereign in some respects and dependent in others:
a stack with European data residency may still be operationally captive
if only the supplier's certified engineers can operate it, and a stack
that is technically open may still be financially captive if the only
realistic supplier holds a monopolistic position. The five dimensions
below are intentionally orthogonal: they ask different questions, can
have different answers for the same component, and require different
remediations.

\subsection{Data sovereignty}
\label{sec:sg-data}

The capacity to determine where institutional data is stored, under
which jurisdiction it is processed, and which actors --- including
state actors --- can compel its disclosure. Data sovereignty depends on
storage location, encryption regime, key custody, and the legal basis
under which a provider may transfer data across borders. The Court of
Justice of the European Union, in \emph{Schrems~II}~\cite{schrems2020},
clarified that contractual safeguards alone do not satisfy this
dimension when data falls under a third-country regime that authorises
disproportionate state access. The European Data Protection Board has
since elaborated the conditions under which transfers remain
permissible~\cite{edpb-guidelines}, conditions that operationalise the
data dimension for any university subject to the
\gls{gdpr}~\cite{gdpr}.

The analytical structure of the data dimension transposes across
jurisdictions: the questions of \emph{where} data sits, \emph{under
which law}, and \emph{accessible to which actors} are universal, even
when the binding instruments differ.

\subsection{Technical sovereignty}
\label{sec:sg-technical}

The capacity to inspect, modify, repair and extend the components of
the infrastructure. Technical sovereignty requires source-code access
where applicable, open data formats, documented application
programming interfaces, and the absence of binary-only critical paths. A component whose internals the
institution cannot examine is, by construction, opaque to its own
governance: incidents cannot be analysed beyond the supplier's
willingness, regressions cannot be confirmed independently, and the
institution's defensive posture is bounded by the supplier's pace.

Technical sovereignty does not require that every component be open
source. It requires that the institution can either inspect a component
itself or substitute it with another that it can. Replaceability is
the operational form of technical sovereignty when inspection is
unavailable.

\subsection{Financial sovereignty}
\label{sec:sg-financial}

The capacity to predict and contain the cost of operating the
infrastructure over time, and to be free of captive escalation.
Financial sovereignty is reduced when pricing depends on a single
supplier with no realistic substitute, when licence audits create
unbounded liability, or when exit costs dominate exit decisions to the
point where no exit is ever attempted. It is restored by genuine
multi-supplier capability, by predictable cost models with contractual
caps, and by exit costs proportionate to the genuine migration work
rather than to the supplier's commercial preference.

Financial sovereignty is the dimension most often treated as a
secondary concern at procurement and as the primary concern at the
next renewal. Treating it as a first-order dimension reverses the
order: it is examined before, not after, the technical evaluation.

\subsection{Operational sovereignty}
\label{sec:sg-operational}

The capacity to operate the infrastructure with skills available
locally or through trusted partners --- the institution's own staff,
its national research and education network, sector consortia, peer
institutions. A component that requires a specific supplier's certified
engineers, available only through that supplier, places the institution
in operational dependence even when other dimensions are satisfied.
Conversely, components built on widely-known standards admit a deep
labour market, peer support, and the possibility of mutualised operation
across institutions.

Operational sovereignty interacts strongly with financial sovereignty:
a high cost of operation may be acceptable when paid to the local
labour market, and unacceptable when paid as a captive engagement to
the same supplier that holds the licence.

\subsection{Institutional sovereignty}
\label{sec:sg-institutional}

The capacity to retain authority over the decisions that matter ---
what to do, when, with whom --- without delegating them to a contractual
arrangement that displaces them. Institutional sovereignty erodes
silently: it is rarely transferred in a single decision but accumulated
over many small ones. It manifests as the inability to refuse a
roadmap change, to choose a non-default integration, to deviate from
a supplier's reference architecture without penalty, or to set internal
priorities that the supplier has not anticipated.

Cybersecurity authorities have emphasised the governance dimension of
digital risk management as a distinct concern from the technical one
--- as the European Union Agency for Cybersecurity has
done~\cite{enisa-risk} --- and the applicable cybersecurity regulatory
regime commonly codifies this distinction by requiring management
bodies to retain accountability for security measures, as the European
NIS2 Directive does~\cite{nis2}. In a university, institutional
sovereignty is the precondition of academic self-government over
digital matters.

%-----------------------------------------------------------------------------
\section{Indicators per dimension}
\label{sec:sg-indicators}

Each dimension is grounded in four measurable indicators. The
indicators are deliberately conservative: they ask whether a specific
capacity \emph{exists}, not whether the institution \emph{exercises it
optimally}. The full canonical table is given in
Table~\ref{tab:sg-indicators}.

\begin{table}[H]
\centering\small
\caption{The twenty indicators of the sovereignty grid (canonical
form). Indicators are ordered by dimension, then by criticality.}
\label{tab:sg-indicators}
\begin{tabularx}{\textwidth}{%
  >{\hsize=0.60\hsize\linewidth=\hsize\bfseries\raggedright\arraybackslash}X%
  >{\hsize=0.90\hsize\linewidth=\hsize\raggedright\arraybackslash}X%
  >{\hsize=1.50\hsize\linewidth=\hsize\raggedright\arraybackslash}X%
}
\toprule
Dimension & Indicator & Question (verifiable capacity) \\
\midrule
\multirow{4}{*}{\textbf{Data}}
  & Storage location & Is the storage location of all classes of data
    documented and contractually fixed? \\
  & Key custody & Are encryption keys for data-at-rest held in custody
    that the institution controls (BYOK or HYOK)? \\
  & Disclosure regime & Has the disclosure regime applicable to the
    data been characterised, including cross-border state access? \\
  & Exit-data export & Can data be exported in a documented open
    format within an agreed time window without supplier
    cooperation? \\
\addlinespace
\multirow{4}{*}{\textbf{Technical}}
  & Source visibility & Is the source code (or equivalent
    specification) available for inspection by the institution or its
    designated experts? \\
  & Standard interfaces & Does the component speak the relevant
    interface standards without captive extensions in critical paths? \\
  & Independent build & Can the operational artifact be built by a
    party other than the supplier from publicly available material? \\
  & Regression auditability & Can a behavioural regression be
    confirmed independently of supplier diagnostics? \\
\addlinespace
\multirow{4}{*}{\textbf{Financial}}
  & Cost predictability & Is the multi-year operating cost predictable
    within a documented envelope? \\
  & Supplier substitutability & Are at least two realistic suppliers
    available for the same functional role? \\
  & Audit liability & Is the supplier's audit-and-true-up clause
    bounded in scope and consequence? \\
  & Exit-cost proportionality & Is the contractual exit cost
    proportionate to the genuine migration work? \\
\addlinespace
\multirow{4}{*}{\textbf{Operational}}
  & Local skills availability & Are the skills required to operate the
    component available through the institution, the \gls{nren} or peer
    institutions? \\
  & Labour-market depth & Does the underlying technology admit a
    labour market broader than the supplier's certified
    pool? \\
  & Documentation completeness & Is the component documented to a
    level sufficient for an independent operator to take
    over? \\
  & Mutualisation feasibility & Can the operational burden be shared
    with peer institutions through an existing or feasible
    consortium? \\
\addlinespace
\multirow{4}{*}{\textbf{Institutional}}
  & Roadmap influence & Can the institution decline a supplier
    roadmap change without contractual penalty? \\
  & Integration freedom & Can the institution adopt non-default
    integrations without operational degradation? \\
  & Decision retention & Are the decisions of consequence (data
    classification, access, retention, incident handling) made by the
    institution rather than imposed by the supplier? \\
  & Public accountability & Can the institution publicly account for
    the configuration in place without violating supplier
    confidentiality clauses? \\
\bottomrule
\end{tabularx}
\end{table}

The indicators are not independent: financial and operational
sovereignty in particular tend to interact, since a position of
supplier dependence on one dimension typically constrains the
institution's ability to recover the other. This is offered as an
analytical observation rather than an empirical measurement, and
should be tested against the institution's own history. Reporting the
dimensions separately preserves analytical clarity, but deliberation
should attend to these interactions rather than treat each dimension
as fully isolated.

%-----------------------------------------------------------------------------
\section{Scoring guidance}
\label{sec:sg-scoring}

Each indicator is scored on a fixed four-point scale:

\begin{description}[leftmargin=2.5em,nosep,style=nextline]
  \item[0 --- Absent] The capacity does not exist. The indicator is
  unsatisfied; remediation requires a structural change.
  \item[1 --- Partial] The capacity exists but with significant
  restrictions (scope, jurisdiction, time window, supplier cooperation
  required).
  \item[2 --- Established] The capacity is present in normal operating
  conditions, documented, and known to the relevant team.
  \item[3 --- Robust] The capacity has been verified by exercise (a
  drill, an audit, or a real event) within the last twenty-four
  months (this Reference adopts twenty-four months by convention, to
  be aligned with the institution's own audit cadence), beyond which
  operational drift renders an unverified capacity unreliable.
\end{description}

Half-points are not allowed. When uncertainty exists, the lower score
is recorded, with the basis for uncertainty noted.

\paragraph{Aggregation.} The grid does \emph{not} produce a single
sovereignty score. Each dimension is reported independently. Two
practical conventions apply:

\begin{itemize}[leftmargin=1.5em,nosep]
  \item \textbf{Per-dimension reporting} --- the dimension score is the
    minimum of its indicator scores. A dimension with a single Absent
    indicator is reported as Absent regardless of the others. This
    reflects the fact that a single critical gap defines the
    institution's exposure on that dimension.
  \item \textbf{Procurement gates} --- per workload class, the
    institution may set lexicographic gates (e.g.\ \emph{data dimension
    must be at least Established for sensitive workloads}). Gates are
    declared in advance, not adjusted to fit the candidate.
\end{itemize}

\paragraph{Cadence.} The grid is exercised at three moments: at
procurement (gate decision), at annual operational review (drift
detection), and at every renewal (re-qualification). Between these
moments, an event-driven re-exercise is triggered by any of: a
significant pricing change, a supplier acquisition, a court decision
or regulatory change affecting the disclosure regime, a major outage,
or a publicly disclosed supplier security incident.

%-----------------------------------------------------------------------------
\section{Worked examples}
\label{sec:sg-examples}

The three examples below score the same functional layer --- identity
provision --- with three different instantiations. The intent is
illustrative, not prescriptive: the same instantiation in a different
institution may receive different scores because the verifiable
capacities are institution-specific.

\subsection{Case 1 --- Pure-OSS (Keycloak + SCIM + eduGAIN)}
\label{sec:sg-ex-oss}

A self-hosted Keycloak deployment, integrated with the institution's
\gls{sso} federation through eduGAIN, with \gls{scim} provisioning to
downstream applications.

\begin{table}[H]
\centering\small
\caption{Sovereignty profile of the pure-OSS identity instantiation.}
\label{tab:sg-ex-oss}
\begin{tabularx}{\textwidth}{%
  >{\hsize=0.75\hsize\linewidth=\hsize\bfseries\raggedright\arraybackslash}X%
  >{\hsize=0.75\hsize\linewidth=\hsize\raggedright\arraybackslash}X%
  >{\hsize=1.50\hsize\linewidth=\hsize\raggedright\arraybackslash}X%
}
\toprule
Dimension & Score & Comment \\
\midrule
Data          & 3 Robust       & Storage under the institution's own data-custody and sovereignty profile, institutional key custody, no third-country disclosure regime. \\
Technical     & 3 Robust       & Source open, standards respected, regressions auditable. \\
Financial     & 2 Established  & Operating cost predictable, no licence; substitutability via Shibboleth or Authentik. \\
Operational   & 2 Established  & Skills available; \gls{nren} community support; no supplier-only certification. \\
Institutional & 3 Robust       & Roadmap controlled internally; integration freedom complete. \\
\bottomrule
\end{tabularx}
\end{table}

\noindent\emph{Profile.} Strong on data, technical and institutional;
operational and financial dimensions are good but limited by the
internal staffing burden. The principal risk is operational fragility
if the maintaining team shrinks below critical mass.

\subsection{Case 2 --- Pure commercial (a dominant commercial identity suite, e.g.\ Microsoft Entra~ID, with conditional-access policies)}
\label{sec:sg-ex-entra}

A managed identity tenancy from a dominant commercial suite vendor,
operated by the supplier in a designated in-jurisdiction region and
integrated through the institution's existing tenant arrangements.

\noindent\emph{Configuration assumed.} The scoring below reflects the
default supported configuration with in-jurisdiction data residency and
standard supplier-managed key custody. Optional arrangements such as
restricted-access lockbox features, double-key encryption, or
\gls{hyok} integrations would modify the data dimension upward, but at
the cost of additional licence and operational commitments; their
adoption is itself a sovereignty-affecting choice that an institution
should evaluate explicitly.

\begin{table}[H]
\centering\small
\caption{Sovereignty profile of the pure-commercial identity instantiation
(default configuration, without optional HYOK arrangements).}
\label{tab:sg-ex-entra}
\begin{tabularx}{\textwidth}{%
  >{\hsize=0.75\hsize\linewidth=\hsize\bfseries\raggedright\arraybackslash}X%
  >{\hsize=0.75\hsize\linewidth=\hsize\raggedright\arraybackslash}X%
  >{\hsize=1.50\hsize\linewidth=\hsize\raggedright\arraybackslash}X%
}
\toprule
Dimension & Score & Comment \\
\midrule
Data          & 1 Partial      & In-jurisdiction regions designated; key custody bounded; disclosure regime under the applicable cross-border transfer analysis remains contested in the absence of \gls{hyok}. \\
Technical     & 1 Partial      & Standards spoken with captive extensions in conditional-access policy expressions; regressions traceable only via supplier diagnostics. \\
Financial     & 1 Partial      & Per-seat pricing with documented historical escalation; substitutability constrained by ecosystem coupling. \\
Operational   & 3 Robust       & Skills market deep; documentation extensive; enterprise support widely available. \\
Institutional & 1 Partial      & Roadmap supplier-driven; integration freedom limited to certified partner products. \\
\bottomrule
\end{tabularx}
\end{table}

\noindent\emph{Profile.} Strong only on operational. The pattern is
characteristic of mature commercial offerings: the dimensions where
the institution's bargaining position is structurally weaker
(technical, financial, institutional) score lower in the default
configuration, while the dimension underwritten by a deep labour
market (operational) is excellent. Neither asymmetry is intrinsic to
commercial supply; both are amenable to contractual remediation when
the institution makes them an explicit gate at procurement.

\subsection{Case 3 --- Hybrid (Shibboleth federation + Keycloak local activation)}
\label{sec:sg-ex-hybrid}

The existing Shibboleth federation is preserved for inter-institutional
authentication; a self-hosted Keycloak instance handles role
activation, attribute enrichment and \gls{tac} session management
locally.

\begin{table}[H]
\centering\small
\caption{Sovereignty profile of the hybrid identity instantiation.}
\label{tab:sg-ex-hybrid}
\begin{tabularx}{\textwidth}{%
  >{\hsize=0.75\hsize\linewidth=\hsize\bfseries\raggedright\arraybackslash}X%
  >{\hsize=0.75\hsize\linewidth=\hsize\raggedright\arraybackslash}X%
  >{\hsize=1.50\hsize\linewidth=\hsize\raggedright\arraybackslash}X%
}
\toprule
Dimension & Score & Comment \\
\midrule
Data          & 3 Robust       & Federation metadata public; activation data held locally. \\
Technical     & 3 Robust       & Both components open; standards respected end-to-end. \\
Financial     & 2 Established  & Two operating teams (federation, local); cost predictable but doubled. \\
Operational   & 2 Established  & Two skill sets required; mutualisation feasible through the \gls{nren}. \\
Institutional & 3 Robust       & Roadmap fully retained; integration freedom complete. \\
\bottomrule
\end{tabularx}
\end{table}

\noindent\emph{Profile.} The hybrid profile illustrates a recurring
pattern: better than pure-commercial on every dimension, comparable to
pure-OSS on most, with a financial and operational cost premium that
reflects the dual stack. The premium is the price of preserving an
existing federation investment without surrendering institutional or
technical sovereignty.

%-----------------------------------------------------------------------------
\section{Common misuses of the grid}
\label{sec:sg-misuses}

Three anti-patterns recur in practice and each defeats the purpose of
the instrument.

\paragraph{Retro-fitting the grid to a pre-decided choice.} A
procurement decision is taken on other grounds (vendor relationship,
ease, urgency) and the grid is then completed to legitimise the
choice. The mitigation is procedural: gates are declared before the
candidate is scored, not after.

\paragraph{Treating the dimensions as fully independent.} The
dimensions are orthogonal in the sense that they ask different
questions, but their answers correlate. A monopoly position
(financial deficit) is almost always accompanied by an inability to
refuse roadmap changes (institutional deficit). Treating them as fully
independent permits the false comfort of compensating one with the
other.

\paragraph{Applying the grid only at procurement.} The sovereignty
profile drifts across the operational life of a component: tariffs
change, suppliers consolidate, regulatory frameworks evolve. A grid
applied only at procurement captures the worst predictive moment ---
the moment of maximum supplier accommodation --- and never re-examines
the trajectory.

%-----------------------------------------------------------------------------
\section{The grid as a governance instrument}
\label{sec:sg-governance}

The grid is not solely a technical artifact. It is the connective
tissue between four governance practices that are frequently conducted in mutual
ignorance.

\begin{description}[leftmargin=1.5em,style=nextline,nosep]
  \item[Procurement.] The grid translates into \gls{rfp} clauses,
    contractual exit assistance, and gating thresholds per workload
    class. It supplies the procurement team with a vocabulary
    independent of supplier marketing categories.

  \item[Internal audit.] The grid provides the structure for an annual
    review of the institution's sovereignty profile, comparable across
    components and across years.

  \item[Peer-institution benchmarking.] When a network of universities
    adopts the same grid, the comparison of profiles ceases to be
    anecdotal. The phase-2 international university network anticipated
    by this Reference is designed around this premise.

  \item[Public accountability.] An institution that publishes its
    sovereignty profile by component renders its digital strategy
    legible to its academic community, its students, and the public
    that funds it. The grid makes this publication concise and
    comparable.
\end{description}

\medskip

\noindent The grid alone does not produce sovereign infrastructure. It
produces the precondition for sovereign infrastructure: a vocabulary
shared across procurement, operations, audit and governance, in which
sovereignty can be discussed without reduction to ideology and without
indefinite deferral.
